\chapter{La población y la vivienda vacía}


\label{cap:poblacion}

Entre 2010 y 2022 el departamento La Caldera creció más que casi cualquier otro distrito
del país, y al final de ese ciclo tenía una de cada cinco viviendas sin nadie viviendo en
ella. Este capítulo mide las dos cosas, muestra que la segunda es anterior a la primera y
llega a una conclusión que ninguna de las dos anticipa: la vacancia y el turismo no son
dos problemas sino dos nombres del mismo uso del suelo.

\section{La magnitud}

El Censo 2022 fija la escala del proceso y conviene enunciarla antes que cualquier interpretación.

\begin{ficha}{Censo Nacional 2022 --- departamento La Caldera}
\begin{itemize}[leftmargin=*,nosep]
\item Población: \textbf{12.299 habitantes} (resultados definitivos). Crecimiento respecto de 2010: \textbf{58,4\%}.
\item Es el \textbf{cuarto departamento con mayor crecimiento demográfico de la Argentina} en la década. Lo anteceden Añelo, Neuquén (68,4\%, de 10.786 a 18.166 habitantes); San Vicente, provincia de Buenos Aires (65,1\%, de 59.478 a 98.215); y General Rodríguez, Buenos Aires (63,7\%, de 87.185 a 142.709).
\item Cerrillos, el otro departamento metropolitano salteño, creció 56,3\%, de 35.789 a 55.949 habitantes.
\item Provincia de Salta: crecimiento del 18,7\%, de 1.214.441 a 1.441.351 habitantes. Departamento Capital --- que incluye los municipios de Salta y San Lorenzo ---: 17,1\%, de 536.113 a 627.704.
\item \textit{Fuente: INDEC, Censo 2022, resultados definitivos, cuadro 1.17 (Salta); el orden nacional, según} Cuarto, \textit{6/2/2023.}
\end{itemize}
\end{ficha}

\begin{cautela}
Una advertencia de fecha. Las cifras salteñas de esta ficha son las \textbf{definitivas} del cuadro 1.17 del INDEC (noviembre de 2023). Una versión anterior usaba las provisorias de enero de 2023 ---1.440.672 para la provincia, 627.107 para la Capital, 56.287 para Cerrillos y 12.458 para La Caldera---, y con ellas daba un 18,6 por ciento provincial, un 57,3 para Cerrillos y a éste el quinto puesto del país. Las definitivas bajan a Cerrillos al 56,3 y dejan a La Caldera, con 58,4, como el departamento que más creció en Salta. El \textbf{orden nacional} está verificado sólo en parte. Los cuadros definitivos de Buenos Aires confirman que San Vicente (65,1 por ciento) y General Rodríguez (63,7) superan a La Caldera, con porcentajes algo menores que los provisorios de la nota periodística (66,4 y 64,3), y que ningún otro partido bonaerense la alcanza: el siguiente, Mar Chiquita, creció 55,6. El cuadro definitivo de Neuquén sube a Añelo del 65,9 provisorio al 68,4, y lo deja primero; ningún otro departamento neuquino supera a La Caldera (el siguiente, Huiliches, creció 42,4). Los tres que la anteceden quedan confirmados. Y los cuadros definitivos de otras siete provincias ---las de mayor expansión demográfica o turística de la década--- no muestran ningún departamento por encima del 58,4: los más altos son Río Grande, en Tierra del Fuego (50,3), Santa María, en Córdoba (47,5), Luján de Cuyo, en Mendoza (46,0), Lago Buenos Aires, en Santa Cruz (44,1), Belgrano, en San Luis (40,1), Cushamen, en Chubut (34,8), y El Cuy, en Río Negro (31,8). Los cuadros de otras trece provincias tampoco: el departamento que más creció en ellas es San Antonio, en Jujuy, con 50,7, seguido de Tafí del Valle, en Tucumán, con 50,3, y de Toay, en La Pampa, con 45,3. El cuarto puesto queda así verificado en veintitrés de las veinticuatro jurisdicciones, con los resultados definitivos de 512 departamentos y partidos con dato comparable entre 2010 y 2022. Sólo falta la Ciudad de Buenos Aires, dividida en comunas. \pendiente{cuadro 1.x definitivo del Censo 2022 de la Ciudad de Buenos Aires}
\end{cautela}

\begin{cautela}
Los tres departamentos que crecieron más que La Caldera describen, entre los tres, el mapa de la expansión argentina de la década. San Vicente y General Rodríguez son el frente de urbanización del conurbano bonaerense. Añelo es Vaca Muerta.

La Caldera creció al ritmo de un yacimiento de hidrocarburos no convencionales y del borde del área metropolitana más grande del país, con una diferencia decisiva: en Añelo el crecimiento tiene una base productiva que financia infraestructura, y en el conurbano bonaerense hay una escala de coparticipación y de aparato estatal que no existe aquí. En La Caldera el crecimiento es enteramente residencial y de absorción metropolitana, sobre un municipio que dos años antes de ese censo declaró que la coparticipación no le alcanzaba para pagar sueldos.

Y creció más del triple que la propia capital provincial, que es de donde viene esa población.

Ese número, además, convierte en una advertencia cumplida el diagnóstico que el propio PDUA había incorporado en 2015 citando el Plan de Ejecución Metropolitano de 2012. Aquel documento consignaba a La Caldera con 37,3 por ciento de crecimiento intercensal, ya el mayor del área, y advertía que sin intervención previa los municipios del conurbano salteño se convertirían en periferia degradada o en base de asentamiento de sectores cerrados. En la década siguiente el crecimiento se aceleró al 58,4 por ciento.
\end{cautela}

La discrepancia entre las dos cifras que circulaban sobre la población departamental quedó resuelta al procesar el cuadro 1.17 del Censo 2022, y la explicación es sencilla: eran dos magnitudes distintas y una serie provisoria.

\begin{ficha}{Censo 2022, resultados definitivos --- departamento La Caldera}
\begin{center}\small
\begin{tabular}{@{}lrrr@{}}
\toprule
& \textbf{La Caldera} & \textbf{Vaqueros} & \textbf{Departamento} \\
\midrule
Viviendas & 2.025 & 2.976 & 5.001 \\
Viviendas particulares & 2.023 & 2.968 & 4.991 \\
\textbf{Población total} & \textbf{4.525} & \textbf{7.774} & \textbf{12.299} \\
Población en viviendas particulares & 4.511 & 7.735 & 12.246 \\
Población en viviendas colectivas & 14 & 39 & 53 \\
Población en situación de calle & --- & --- & --- \\
\bottomrule
\end{tabular}
\end{center}
\textit{Fuente: INDEC, Censo 2022, cuadro 1.17, total de viviendas y población según gobierno local. Categoría de ambos municipios: M3.}
\end{ficha}

\begin{cautela}
La cifra de \textbf{12.246} que este libro venía usando como ``población en hogares'' es correcta y es definitiva: es la población en viviendas particulares del departamento, suma exacta de los dos municipios. La cifra de \textbf{12.458} pertenece a los \textbf{resultados provisorios} difundidos en febrero de 2023 y no debe usarse. El total definitivo es \textbf{12.299}, que es 12.246 más las 53 personas censadas en viviendas colectivas.

En consecuencia, el crecimiento intercensal del departamento es del \textbf{58,4 por ciento} --- de 7.763 a 12.299 entre 2010 y 2022 --- y no del 60,5 que daban las cifras provisorias. Sigue siendo \textbf{3,1 veces} el crecimiento provincial del 18,7 por ciento, de modo que ninguna afirmación del libro cambia de sentido; cambian los decimales, y conviene que sean los correctos.

El departamento está íntegramente cubierto por los dos municipios: no hay área sin gobierno local. \textbf{Vaqueros concentra el 63 por ciento de la población} departamental y el 59 por ciento de las viviendas particulares.
\end{cautela}

\subsection*{Dos personas y cuarto por vivienda}

Procesado el cuadro para los sesenta gobiernos locales de la provincia, aparece un indicador que corrobora por otra vía el hallazgo de vacancia de este libro.

{\small
\begin{longtable}{@{}lrr@{}}
\toprule
& \textbf{Personas por vivienda particular} & \textbf{Puesto (de 60)} \\
\midrule
\endhead
Provincia de Salta & 3,13 & --- \\
Vaqueros & 2,61 & 14º más bajo \\
Departamento La Caldera & 2,45 & --- \\
\textbf{Municipio de La Caldera} & \textbf{2,23} & \textbf{3º más bajo} \\
\bottomrule
\end{longtable}}

El indicador hay que leerlo con cuidado, porque tiene dos causas posibles y sólo una interesa a este libro.

Los dos municipios que anteceden a La Caldera en el ranking son \textbf{Los Toldos} (1,87) y \textbf{Nazareno} (2,10), y la lista de los diez más bajos se completa con El Potrero, Payogasta, San Carlos, Cachi, Isla de Cañas, Santa Victoria Oeste y Guachipas. Son municipios de la Puna, los Valles y las yungas fronterizas, donde una baja relación entre personas y viviendas expresa \textbf{despoblamiento}: casas que quedaron vacías porque la gente se fue.

La Caldera está en esa lista por la razón contraria. Pertenece al departamento que más creció de la provincia y el cuarto que más creció del país, y tiene la tercera menor ocupación por vivienda de Salta. \textbf{Crecimiento y vacancia simultáneos} no se explican por emigración: se explican porque una parte del parque construido no está destinada a residencia permanente.

Ésa es exactamente la conclusión que este libro había obtenido por vía independiente, contando viviendas sin residentes permanentes en Redatam. Dos fuentes distintas, dos métodos distintos, el mismo resultado. Es la corroboración más sólida que el capítulo tiene.


Un antecedente, antes de la serie histórica: la literatura urbanística describe el corredor San Lorenzo--Vaqueros--La Caldera como zona donde la expansión adopta la forma de emprendimientos residenciales de muy baja densidad, en gran parte cerrados y orientados a estratos medios y altos, que avanzan sobre los ecosistemas de transición de Yungas. Consigna asimismo que entre 1970 y 2010 se perdieron 3.754 hectáreas de suelo productivo sólo en el municipio de Salta.


\section{La casa quinta es de 1913}

Antes de medir la vacancia con los censos conviene establecer desde cuándo existe, porque
este libro sostuvo durante buena parte de su elaboración que era un efecto del ciclo
inmobiliario reciente. No lo es.

\textbf{El primer indicio es de 1913 y son dos avisos del mismo año.} En junio se remata «la
finca Lesser», de mil seiscientas cuadras cuadradas, y el aviso no la vende como campo sino
como paisaje: está «\textit{ubicada en el departamento La Caldera, el punto más pintoresco de
la provincia de Salta}», tiene «\textit{una hermosa casa de campo, grandes rastrojos de
cultivo, quinta, abundante agua clara durante todo el año}» y queda «\textit{a una hora de
viaje en carruaje de la ciudad}»; el remate se hace en la confitería del Jockey Club, en la
plaza 9 de Julio de la capital (B.O.\ Nº 399, 1913, h.\ 4). En julio se remata, en la capital
también, \textbf{una casa-quinta en el pueblo mismo}, «\textit{de siete espaciosas habitaciones
y sobre una cuadra cuadrada de terreno}», con base de \$1.500 y con don Jacinto
Guerrero\index{Guerrero, Jacinto} de lindero por el poniente, el sur y el norte (Nº 409, h.\ 3).

\textbf{El vendedor, en 1913, ya le habla a un comprador de la ciudad, y le vende clima,
paisaje y distancia.} Es lo mismo que sostiene el relato de los vecinos que este capítulo
registra más abajo, catorce años antes que el primer remate de casa quinta que este libro
conocía.

\begin{ficha}{Remate judicial --- B.O.\ Nº 1.192, 11 de noviembre de 1927}
Por orden del juez C.\ Gómez Rincón, y sobre bienes embargados por el \textbf{Banco de la
Nación Argentina} en el juicio ejecutivo 13906 de 1927, el martillero Antonio Forcada
remata:

\textbf{«Una casa quinta ubicada en Lesser, Departamento de La Caldera, de 10 hectáreas más
o menos, compuesta de tres fracciones unidas.»} Base: \$6.000 al contado.

Los límites de las tres fracciones se describen por accidentes: «una quebradita situada en
la cumbre del camino a Lesser»; «cumbre del cerro del Potrerillo»; «propiedad de Zerda»;
«finca Lesser»; «río Lesser»; «la última barranca del río Castellanos».

En el mismo aviso se rematan los derechos indivisos de una casa en calle 20 de Febrero de
la capital, tasada en \$70.000.
\end{ficha}

Una \textbf{casa quinta} no es una explotación agropecuaria ni una vivienda permanente: es
una residencia de descanso. Diez hectáreas en Lesser, con acceso por camino propio y frente
sobre el río, embargadas a alguien que además tenía derechos sobre una casa en el centro de
Salta.

El uso recreativo del suelo de este departamento por parte de propietarios que viven en la
capital está documentado en 1913 y vuelve a estarlo en 1927. Cuando el capítulo mida que una de cada cinco viviendas no
tiene residentes permanentes, y que esa proporción ya era del diecinueve por ciento en 2010
—es decir, anterior al ciclo de loteos que este libro documenta—, conviene tener presente
que la serie es todavía más larga: la segunda residencia caldereña tiene por lo menos un
siglo.

Y no fue un caso aislado. Dos años después, en octubre de 1929, se remató por orden
judicial \textbf{otra casa quinta, esta vez en Vaqueros}, «con una extensión de cinco
hectáreas más o menos», con base de mil pesos —«las dos terceras partes de su avaluación»—
y con estos límites: «al Norte, con la propiedad de la Señora de Bianchi; Sud con el
\textbf{Río Vaqueros}; Este con propiedad de Arturo Fernández, y Oeste con la finca del
Doctor Carlos Outes».

Dos casas quinta rematadas en dos años, una en Lesser y otra en Vaqueros, ambas con frente
sobre un río, ambas de entre cinco y diez hectáreas, ambas ejecutadas por deudas. Y los
linderos son un abogado de la capital ---el doctor Carlos Outes--- y dos apellidos más:
\textbf{Fernández}, que el archivo del departamento registra en el comisario auxiliar de San
Alejo y Santa Rufina de 1909, en el comprador de la finca Campo Alegre de 1937 y en el
presidente de la Comisión Municipal de 1983; y \textbf{Bianchi}, que el capítulo
\ref{cap:politica} encuentra en 1928 denunciando a Urquiza por uso indebido de agua.
Coincidencias de apellido, como todas las del apéndice \ref{cap:personas}.

Lo que el ciclo reciente hizo no fue inventar ese uso. Fue multiplicarlo por escala y
formalizarlo en loteos.

\subsection*{Una hipótesis que el archivo permite formular, y sus límites}

Es tentador dar un paso más y decir que el departamento \textit{fue pensado} como lugar de
residencia de verano para la clase alta salteña. El archivo da elementos para sostenerlo y da
otros para resistirlo, y conviene ver los dos antes de elegir.

\textbf{Lo que lo sostiene no es la casa quinta: son las truchas.} En julio de 1940 la
Provincia autoriza \$673,45 para traer al Jefe de la División de Piscicultura de la Nación y
comprar \textbf{diez mil truchas}; en noviembre de 1942 paga \$200 por \textbf{cinco mil
alevinos de trucha arco iris sembrados en los ríos Castellanos y Wierna} (capítulo
\ref{cap:infraestructura}). Sembrar trucha en un río no alimenta a nadie: es infraestructura
de pesca deportiva. El informe de 1940 nombra además, entre las aguas privadas ya pobladas de
pejerrey, la finca de \textbf{Lucio Ortiz} en Yacones ---el mismo a quien en 1942 se le
reconoce derecho exclusivo sobre un arroyo entero, capítulo \ref{cap:redes}--- como una de las
cuatro instalaciones disponibles de la provincia.

Póngase eso junto a lo demás. En 1935 la Legislatura vota \textbf{\$7.000 para ampliar la
Iglesia Parroquial}, diez veces lo liquidado ese mismo año para impedir que el río entrara a
las casas. En el revalúo de 1944, \textbf{Elvira Patrón Costa de Solá y Elina Solá Patrón de
Ortiz} ---los apellidos que gobernaron Salta esa década--- suman casi el diez por ciento del
valor del departamento, y \textbf{Carlos Serrey}\index{Serrey, Carlos} tiene casa en el casco y
rodea a la Curia por dos costados (capítulo \ref{cap:fincas}). Trucha sembrada, iglesia
ampliada, apellidos de gobierno entre los mayores contribuyentes y agua privada declarada
exclusiva por decreto.

\textbf{Y lo que lo resiste es que el Estado hizo lo contrario, dos veces por escrito y una
por omisión.} En 1934 la Ley 1423 (orig.\ 145) incorpora a La Caldera a la \textbf{región económica del
olivo} y exime por diez años de todo impuesto a quien plante olivos; en 1980 el Decreto 12
reconoce su \textbf{área tabacalera} (capítulos \ref{cap:hacienda} y \ref{cap:tierra}). Si el
departamento hubiera estado pensado como veraneo, esos dos actos no existirían. Existen, y
ninguno de los dos perfiles sobrevivió. Y el gasto público tampoco acompaña a un destino: en
1935 los vecinos piden un puente y la respuesta es «en su oportunidad»; en \textbf{mayo de 1944
se elimina el camino de Vaqueros a Los Yacones del plan de obras} y sus \$40.000 se reasignan
fuera del departamento porque «el gasto no estaría justificado con el beneficio que la
inversión reportaría»; en 1937 se niega la aquiescencia para una escuela en Potrero de
Castilla ---y se concede veinticuatro días después, sólo porque la Inspección Nacional
comprobó que no había escuela provincial en el lugar---. El capítulo \ref{cap:presencia} lo resume: entre 1923 y 1945 el Estado puso policía,
registro, recaudación y vías de paso, y los caminos que sí construyó van a Campo Santo, a
Güemes, a Jujuy. \textbf{De paso, no de destino.}

\begin{cautela}
De modo que la hipótesis se sostiene en una versión y no en la otra, y la diferencia entre
ambas no es de matiz.

\textbf{El uso residencial-recreativo no fue asignado al departamento: fue el uso más antiguo
que se instaló en él, anterior a los dos perfiles productivos que el Estado intentó imponerle y
que fracasaron. La clase alta salteña no encontró un lugar diseñado para veranear: encontró un
lugar barato, cercano, con agua y con río, y veraneó ahí mientras el Estado lo clasificaba en
tercera categoría y le sacaba los caminos del plan.}

Esa formulación explica tres cosas que la versión fuerte no explica. Explica por qué la
vacancia ya era del diecinueve por ciento en 2010 y por qué la serie arranca en 1927: no es un
efecto del ciclo de loteos, es la forma antigua. Explica el dato que hoy parece anómalo y no lo
es ---que hoteles y restaurantes ocupen a doscientas dieciocho personas, el 3,49 por ciento,
\textit{por debajo} del promedio provincial (capítulo \ref{cap:trabajo})---: el veraneo
caldereño nunca fue una economía turística, fue casa propia de gente de afuera, que es
exactamente lo que el censo mide como vivienda sin residentes. Y no rompe el capítulo
\ref{cap:prospectiva}.

\textbf{Ese último punto es el que obliga a la precisión.} La tercera condición del dispositivo,
en el sentido estricto que este libro adopta, es \textbf{que no haya estratega}; y aquel
capítulo, que el lector encontrará más adelante, enuncia como primera condición de refutación que aparezca un acto de gobierno que
asigne al corredor un perfil que se haya sostenido. \textbf{Si La Caldera fue \textit{pensada}
como veraneo de la elite, hay estratega y el término del título no corresponde.} Si fue
\textit{usada} como veraneo mientras el Estado le probaba vocaciones que no prendían, el
dispositivo queda en pie y gana profundidad: el mercado inmobiliario de los años dos mil no
inventó el uso ni lo trajo de afuera, \textbf{lo escaló}.

\textbf{Y esto no se resuelve razonando: se resuelve con tres documentos}, los tres pedidos en
el apéndice \ref{cap:pedidos}. El \textbf{padrón catastral del revalúo de 1944}, que tiene
cuarenta y siete titulares y diría cuántos domiciliaban en la capital ---el formato existe: los
pedimentos mineros de 1925 consignan «Bernardo Austerlitz, domiciliado en Alberdi 374 de la
capital» (capítulo \ref{cap:siglo})---. Las \textbf{ordenanzas de impuestos del municipio} de
1930, 1943, 1944 y 1945, que el archivo menciona y no publica: si alguna distingue la
\textbf{casa quinta} como categoría tributaria propia, ahí está el Estado nombrando el uso en
lugar de deducirlo de un remate. Y los \textbf{libros parroquiales}, que el capítulo
\ref{cap:ausencias} señala como la serie demográfica más antigua y continua de cualquier pueblo
argentino: un pueblo de veraneo tiene bautismos y matrimonios de gente de la capital celebrados
en verano, y eso es un patrón estacional que se ve a simple vista.

\textbf{Hasta que los tres vuelvan, esto es una hipótesis y no debe citarse como hallazgo de
este libro.}
\end{cautela}

\subsection*{Cuántos votaban, y dónde}

Antes de seguir, dos cifras que el archivo temprano deja y que acotan el tamaño del pueblo. El
decreto de correlación del registro cívico de agosto de 1909 (B.O.\ Nº 89, h.\ 4) implica \textbf{ciento setenta y
siete inscriptos} en el departamento. Tres años después, la ubicación de las mesas para la
elección de electores de gobernador de septiembre de 1912 (B.O.\ Nº 361, h.\ 3) reparte a los empadronados
\textbf{del 1 al 524} en un solo colegio electoral y dos locales: \textbf{la Escuela del Pueblo y
el Juzgado de Paz}. No son censos y no se leen como tales ---son varones inscriptos---, pero
dicen dos cosas: que el padrón se triplicó en tres años y que los dos únicos edificios públicos
que el pueblo podía ofrecer para votar eran la escuela y el juzgado.

\textbf{Y hay una tercera cifra, tres años posterior, que corta esa curva.} Los dos decretos de
ubicación de mesas de 1915 ---el de febrero, para la elección de senadores y diputados del 7 de
marzo, y el de agosto, para la de electores de gobernador--- dan al departamento \textbf{dos mesas} y
no una: la mesa 1 con el padrón \textbf{del 1 al 298}, en la \textbf{Escuela Elemental Mixta del
Pueblo}, y la mesa 2 con el padrón \textbf{del 1 al 216}, en la \textbf{Escuela Infantil Provincial de
Vaqueros} (B.O.\ Nº 541, h.\ 2, y Nº 567, h.\ 2). Sumados, \textbf{quinientos catorce} inscriptos,
tres años después de los quinientos veinticuatro de 1912: \textbf{el padrón dejó de crecer}. Lo que sí
cambió es dónde se vota: el departamento pasó de un colegio electoral con dos locales en el pueblo a
dos mesas separadas por el río, una en el pueblo y otra en Vaqueros.

\begin{cautela}
\textbf{Dos reservas sobre estas cifras, y la primera puede invertir la conclusión.} El primer dígito
del segundo tramo está empastado en la edición de agosto y admite leerse \textbf{216 o 316}; la de
febrero imprime 216 sin duda, y por eso es la cifra que se usa, pero con 316 el total sería
seiscientos catorce y el padrón habría seguido creciendo. \textbf{Las dos lecturas se consignan y no se
corrige ninguna.} La segunda: los tramos de las dos mesas arrancan los dos en 1, de modo que se suman
como dos series y no como una; que no se solapen es lo que el propio decreto supone al partirlas por
local, y no algo que este libro haya verificado.
\end{cautela}

\textbf{Lo que no admite dos lecturas es el edificio.} En 1912 los dos locales de votación eran la
escuela del pueblo y el juzgado de paz, los dos en el pueblo. En 1915 aparece una \textbf{escuela
provincial en Vaqueros}: es la primera que este relevamiento registra en esa banda del departamento, y
es ocho años anterior a la ventana desde la cual el capítulo \ref{cap:presencia} mide la presencia
material del Estado. \textbf{Y dos años después el archivo le pone nombre a quien la dirigía}: el decreto 1.450 del 5 de diciembre de 1917 acuerda la jubilación extraordinaria, «\textit{por motivos de salud}», de \textbf{Sara Blasco}, directora de la Escuela Infantil de Vaqueros, a razón de \$42,24 mensuales, «\textit{el promedio de sueldos en proporción a los años de servicios prestados}» (B.O.\ Nº 681, h.\ 2). El importe dice, de paso, cuánto valía ese cargo: poco más de la mitad de lo que el municipio pagaba en 1908 por limpiar el camino nacional durante todo un año.

\subsection*{Lo que cuentan los vecinos: el río como cerrojo}

Con lo anterior ya escrito, dos vecinos del departamento contaron lo mismo y agregaron algo que
el archivo no decía. Es el primer testimonio oral que este libro incorpora, y se registra como
lo que es: una conversación informal, no una entrevista estructurada, sin grabación. Se los
cita sin nombre, como el libro hace con los particulares que no son parte de un acto público.

\begin{ficha}{Testimonio oral --- dos vecinos de La Caldera, septiembre de 2026 (conversación informal, sin grabación)}
En sustancia, y no en sus palabras: La Caldera era el lugar donde \textbf{políticos salteños y
dueños de fincas y de tierras} pasaban el verano, porque es más fresca que la capital.
\textbf{No había puente}: llegaban en carreta antes de la crecida, y mientras el río venía
crecido \textbf{quedaban aislados del resto de la provincia}, con mucha privacidad. Cuando el
río bajaba, volvían a la ciudad y a su actividad política y laboral.

\textit{No consta si lo que cuentan es recuerdo propio o relato transmitido, ni de quién.}
\pendiente{registro de las conversaciones con dos vecinos sobre el veraneo: fecha, si lo
relatado es recuerdo propio o transmitido y por quién, y consentimiento para citarlas sin nombre}
\end{ficha}

El relato contiene cuatro afirmaciones distintas, y el archivo trata a cada una de un modo
distinto.

\textbf{Que no había puente: el archivo lo confirma.} En agosto de 1935 los vecinos piden a
Vialidad «un puente de unión entre dicho Pueblo y el Camino Nacional Salta-Jujuy», y la
respuesta es que se dispondrá lo que corresponda «en su oportunidad» (capítulo
\ref{cap:defensas}): el pedido sólo tiene sentido si ese puente no existía. En febrero de 1940
un informe de Obras Públicas dice que \textit{se estudiaba} la construcción de un puente, y en
octubre de ese año el plan de quince años de la Ley 12.625 asigna \$200.000 a un Puente Río
Caldera que no vuelve a aparecer en las ediciones relevadas (capítulo
\ref{cap:infraestructura}). Cuándo se habilitó el primero, este libro no lo sabe. Sí sabe que,
noventa años después, el pueblo tiene una sola entrada (capítulo \ref{cap:redes}).

\textbf{Que moverse era difícil: el Estado lo dice en 1911.} Dos decretos de ese año lo ponen
por escrito. En octubre, el Poder Ejecutivo crea una oficina del Registro Civil en la estación
Mojotoro (B.O.\ Nº 291, 1911, h.\ 3) «\textit{atendiendo la larga distancia que lo separa del pueblo de La Caldera}» y
«\textit{las dificultades con que tropieza la gente para trasladarse a aquel pueblo por la falta
de medios de movilidad y de transporte}». Y en diciembre crea una plaza de agente para la
comisaría de Vaqueros (Nº 305, h.\ 4) «\textit{con motivo de la concurrencia de trabajadores en la obra del
puente que se construye sobre el río del mismo nombre}» y en la explotación de canteras junto al
puente sobre el río Mojotoro. \textbf{En 1911 se estaban construyendo los puentes del camino,
pero sobre los ríos Vaqueros y Mojotoro, no sobre el de La Caldera}: veinticuatro años después,
en 1935, el pueblo seguía pidiendo el suyo (capítulo \ref{cap:defensas}). El relato de los
vecinos, en su parte física, tiene ahora el respaldo de la administración que lo dice de sí
misma.

\textbf{Que el río crecía todos los veranos: el archivo lo confirma; que cortara el paso, no.}
En enero de 1938 los propietarios de La Calderilla fundan su pedido de defensas en que
«todos los años en la época de lluvias» el río amenaza sus cultivos (capítulo
\ref{cap:defensas}), y la época de lluvias de Salta es el verano, que es también la estación del
veraneo. Pero ninguno de los actos relevados registra que el pueblo quedara incomunicado, y
el decreto de 1938 habla de cultivos, no de caminos. Es lo que habría que buscar en la prensa de
las temporadas de lluvias. \pendiente{prensa salteña de las temporadas de lluvias, de 1930 a
1960: noticias de La Caldera incomunicada por crecidas y crónica social de veraneantes}

\textbf{Que veraneaban allí políticos y propietarios: el archivo lo hace plausible y no lo
prueba.} Es la hipótesis de las páginas anteriores, con los mismos indicios ---truchas
sembradas, iglesia ampliada, los apellidos que gobernaban entre los mayores contribuyentes de
1944--- y los mismos tres documentos pendientes. El testimonio suma una fuente de otro tipo,
pero no del todo independiente: dos vecinos de un mismo pueblo pueden repetir una misma
tradición local, y lo que un pueblo recuerda de sí mismo también se fija por repetición.

\textbf{Que buscaban el aislamiento por la privacidad: el archivo no puede decirlo.} Un motivo
no queda asentado en un decreto. Esa parte del relato es testimonio y queda como testimonio.

Lo nuevo, entonces, no es el veraneo sino \textbf{el mecanismo}. Si el relato es exacto, la
temporada del veraneo coincidía con la temporada en que el río cerraba el paso, y la falta de
puente, que para el vecino que trabajaba la tierra era un perjuicio, para el propietario que
veraneaba era una condición del lugar. Es la asimetría que este libro encuentra en todo el
departamento, vista desde el agua: \textbf{la misma carencia pesa distinto según quién la
tenga.}

Y conviene decir con la misma claridad lo que el testimonio no autoriza. \textbf{Que a alguien
le conviniera el aislamiento no prueba que alguien lo produjera.} La respuesta dilatoria de
1935 cae dentro de un ajuste presupuestario de Vialidad que la explica sin necesidad de otra
razón, y en 1940 el propio Estado puso el puente en su plan con doscientos mil pesos. Esos dos
documentos van en contra de un aislamiento buscado. Si apareciera uno que mostrara lo
contrario ---una oposición de propietarios al puente, una partida desviada a pedido de
alguien---, habría estratega, y por lo dicho más arriba el término del título dejaría de
corresponder. Hasta ahora no apareció.

\begin{cautela}
Con el testimonio, la hipótesis del veraneo queda así. \textbf{La falta de puente está
documentada, y también las crecidas anuales de la época de lluvias. El veraneo de la elite es
plausible y no está probado. El aislamiento como privacidad buscada es sólo testimonio. Y que la
falta de puente se mantuviera porque convenía no tiene respaldo y tiene dos documentos en
contra.} Es una subhipótesis de la anterior, no una tesis del libro, y tampoco debe citarse como
hallazgo hasta que vuelvan los documentos pedidos.
\end{cautela}




\section{Las viviendas, y una hipótesis que el dato no sostiene}

El Censo permite además medir el parque habitacional, y ese cruce era la prueba que este libro esperaba para su argumento sobre el propietario no residente.

\begin{ficha}{Viviendas particulares --- departamento La Caldera y provincia}
\begin{itemize}[leftmargin=*,nosep]
\item Departamento La Caldera: de \textbf{2.578 viviendas particulares en 2010 a 4.991 en 2022} --- esta última cifra verificada por procesamiento Redatam directo. Crecimiento del \textbf{93,6\%}.
\item Población del departamento: de 7.763 a 12.299 habitantes. Crecimiento del \textbf{58,4\%} según resultados definitivos.
\item Provincia de Salta: de 315.186 a 457.992 viviendas particulares (\textbf{45,3\%}), con un crecimiento de la población en viviendas particulares del \textbf{19,2\%}.
\item Departamento Capital: de 136.743 a 200.640 viviendas. Cerrillos: de 9.243 a 16.820.
\item \textit{Fuentes: informe de gobiernos locales de la Dirección General de Estadísticas de Salta sobre datos del INDEC; cifras departamentales consignadas por} El Tribuno, \textit{27/11/2023.}
\end{itemize}
\end{ficha}

En La Caldera las viviendas crecieron un 93,6 por ciento y la población un 58,4: treinta y cinco puntos de diferencia. Ese contraste, por sí solo, no prueba nada, porque el tamaño medio del hogar cae en todas partes por razones demográficas ajenas al mercado de suelo. En la provincia la relación entre población y viviendas particulares pasó de 3,82 a 3,13 --- una caída del 18,0 por ciento --- y en el departamento La Caldera de 3,01 a 2,46 ---sobre población total---, una caída del 18,2. La \textit{variación} es prácticamente la misma, y quien se detenga ahí concluirá, erróneamente, que no hay nada que explicar.

El nivel, en cambio, es muy distinto, y el recuento de hogares muestra por qué.

\begin{ficha}{Censo 2022 --- área 66077, departamento La Caldera. Composición de hogares}
\begin{itemize}[leftmargin=*,nosep]
\item Población en viviendas particulares: \textbf{12.246} personas. (La población total del departamento es 12.299; la diferencia de 53 corresponde a viviendas colectivas.)
\item Jefas, jefes o personas de referencia del hogar: \textbf{4.003}. Es decir, \textbf{4.003 hogares}.
\item Composición: cónyuge o pareja 2.046 (16,71\%); hijas e hijos, hijastras e hijastros 4.738 (38,69\%); nietas y nietos 656 (5,36\%); nueras y yernos 166; madres y padres 126; otros familiares 414; otros no familiares 51; servicio doméstico y sus familiares 7. (Con los 4.003 jefes, las categorías enumeradas suman 12.207; las 39 personas restantes corresponden a vínculos que el cuadro no desagrega.)
\item \textbf{Personas por hogar: 3,06.}
\item \textbf{Viviendas particulares del departamento: 4.991.}
\item \textit{Fuente: procesamiento Redatam del Censo Nacional 2022, área 66077.}
\end{itemize}
\end{ficha}

El cruce de las dos cifras da el número que este libro venía buscando sin encontrarlo.

Hay \textbf{4.991 viviendas particulares} en el departamento y \textbf{4.003 hogares} que las habitan. La diferencia es de \textbf{988 viviendas particulares sin hogar residente: el 19,8 por ciento del parque}. Una de cada cinco.

Y el dato no se explica por hogares pequeños. El tamaño medio del hogar caldereño es de 3,06 personas, una cifra enteramente ordinaria. Lo que produce la baja relación entre población y viviendas --- 2,45 contra 3,13 provincial --- no es que la gente viva en hogares reducidos: es que hay un quinto del parque habitacional sin gente adentro.

Dicho de otro modo: para alcanzar la relación población/vivienda del promedio provincial, el departamento necesitaría 3.911 viviendas. Tiene 4.991.

\section{La comparación y el desglose}

El procesamiento del Censo por fracción censal permite ahora dos cosas que el apartado anterior dejaba pendientes: comparar el departamento con la provincia y desagregarlo internamente.

\begin{ficha}{Viviendas particulares sin hogar residente, Censo 2022}
\begin{itemize}[leftmargin=*,nosep]
\item \textbf{Provincia de Salta} (159 fracciones censales procesadas): 456.571 viviendas particulares ---1.421 menos que las 457.992 del informe provincial: las tasas de este libro se calculan sobre el procesamiento propio---, 417.533 hogares. Diferencia: 39.038 viviendas sin residentes, el \textbf{8,55\%} del parque. Personas por hogar 3,43; por vivienda 3,13.
\item \textbf{Departamento La Caldera}: 4.991 viviendas, 4.003 hogares. Diferencia: 988, el \textbf{19,80\%}. Personas por hogar 3,06; por vivienda 2,45.
\end{itemize}
\textbf{Por fracción censal:}
\begin{itemize}[leftmargin=*,nosep]
\item Fracción 66077-01: 2.968 viviendas, 2.545 hogares, 7.735 personas. Viviendas sin hogar \textbf{14,25\%}. Personas por hogar 3,04; por vivienda 2,61.
\item Fracción 66077-02: 1.474 viviendas, 1.146 hogares, 3.747 personas. Viviendas sin hogar \textbf{22,25\%}. Personas por hogar 3,27; por vivienda 2,54.
\item Fracción 66077-03: 549 viviendas, 312 hogares, 764 personas. Viviendas sin hogar \textbf{43,17\%}. Personas por hogar 2,45; por vivienda \textbf{1,39}.
\end{itemize}
\textit{Fuente: elaboración propia con base en datos del INDEC, Censo Nacional 2022, procesado con Redatam 7 (CEPAL/CELADE).}
\end{ficha}

La comparación resuelve la pregunta que el apartado anterior dejó abierta. \textbf{La tasa de viviendas sin hogar residente del departamento La Caldera es del 19,80 por ciento; la de la provincia de Salta es del 8,55.} El departamento más que duplica el nivel provincial.

Y el desglose muestra que el fenómeno no está repartido. La fracción 66077-03 tiene un 43,17 por ciento de sus viviendas sin residentes y una relación de 1,39 personas por vivienda: sobre 549 viviendas, hay 764 personas. En un ordenamiento de las 159 fracciones censales de la provincia por tasa de viviendas sin hogar, \textbf{esa fracción ocupa el quinto lugar}, y es la mayor en tamaño de las cinco primeras: las cuatro que la anteceden tienen 53, 54, 147 y 250 viviendas respectivamente, magnitudes en las que un puñado de casos altera el porcentaje. La fracción 66077-02 ocupa el puesto 23 y la 66077-01 el 51.

Es decir: las tres fracciones del departamento están por encima de la media provincial, y una de ellas está entre las cinco con más viviendas sin hogar de Salta.

\begin{cautela}
Corresponde acotar el hallazgo con la misma precisión con que se lo enuncia.

\textbf{Primero, la identificación de las fracciones.} Por magnitud demográfica cabía suponer que la fracción 01, con 7.735 habitantes, es el municipio de Vaqueros, y que las otras dos se reparten el de La Caldera. El capítulo \ref{cap:ausencias} lo confirma con la cartografía censal: la 02 contiene el pueblo y La Calderilla, y la 03 el sector oriental de El Gallinato.

\textbf{Segundo, qué significa ``sin hogar residente''.} La cifra resulta de restar hogares a viviendas particulares, y el Censo cuenta las viviendas particulares habitadas y deshabitadas. Una vivienda sin hogar puede ser una segunda residencia, una unidad en venta o alquiler, una obra en construcción o una vivienda abandonada. La cifra mide la magnitud del fenómeno, no su composición.

\textbf{Tercero, lo que sí se puede afirmar.} Cualquiera sea la composición, el departamento tiene una proporción de vivienda sin residentes que duplica la provincial, y contiene una fracción censal que está entre las cinco más altas de Salta. Eso no se explica por hogares pequeños: el tamaño medio del hogar caldereño es de 3,06 personas, algo más de un diez por ciento por debajo de las 3,43 provinciales, mientras que la relación por vivienda cae de 3,13 a 2,45. La diferencia la produce la vivienda vacía, no la familia chica.
\end{cautela}

Corresponde además una corrección expresa. Una versión anterior de este capítulo sostuvo, sobre la sola comparación de tasas de crecimiento entre 2010 y 2022, que el Censo no aportaba evidencia de un excedente de viviendas y que el argumento del propietario no residente se apoyaba únicamente en la intimación municipal de agosto de 2024. Esa afirmación era incorrecta y su error fue metodológico: comparaba variaciones cuando debía comparar niveles, y no disponía del recuento de hogares.

Con el dato de hogares y la comparación provincial, el argumento tiene ahora respaldo censal. Se deja la corrección visible porque el modo en que un relevamiento se equivoca es información sobre el objeto: la variación del tamaño del hogar oculta, si no se controla, exactamente el fenómeno que se busca medir.

\section{El Censo 2010: la vivienda vacía es anterior al ciclo}

La misma relación puede calcularse hacia atrás, y el resultado corrige la dirección en que este capítulo venía inclinándose.

\begin{ficha}{Comparación intercensal}
\begin{itemize}[leftmargin=*,nosep]
\item \textbf{Departamento La Caldera, 2010:} 2.578 viviendas particulares y \textbf{2.088 hogares}. Diferencia: 490, el \textbf{19,01\%}. Personas por hogar: 3,72.
\item \textbf{Departamento La Caldera, 2022:} 4.991 viviendas y 4.003 hogares. Diferencia: 988, el \textbf{19,80\%}. Personas por hogar: 3,06.
\item \textbf{Provincia de Salta, 2010:} 315.186 viviendas particulares y 299.822 hogares. Diferencia: 15.364, el \textbf{4,87\%}. \textit{(Esta tasa se calcula sobre el recuento completo de viviendas del informe provincial y la de 2022 sobre el procesamiento propio: con la cifra del informe ---457.992--- la de 2022 sería 8,83 en lugar de 8,55. La serie cruza esa diferencia de base, de 1.421 viviendas, y ninguna de las dos lecturas altera su dirección.)}
\item \textbf{Provincia de Salta, 2022:} 456.571 viviendas y 417.533 hogares. Diferencia: 39.038, el \textbf{8,55\%}.
\end{itemize}
\textit{Fuentes: 2022, procesamiento Redatam sobre la base completa. 2010, hogares por procesamiento Redatam sobre la base del cuestionario ampliado con expansor; viviendas particulares totales según el recuento completo consignado en el informe provincial y en} El Tribuno.
\end{ficha}

El resultado es inequívoco en su dirección. \textbf{La tasa de viviendas sin hogar del departamento La Caldera era del 19,01 por ciento en 2010 y es del 19,80 en 2022: prácticamente la misma.} Mientras tanto, la tasa provincial pasó del 4,87 al 8,55, es decir que casi se duplicó.

De modo que la vivienda vacía en este departamento \textbf{no es un producto del ciclo de loteos de la última década}. Ya estaba allí antes, y en 2010 era todavía más anómala en términos relativos: cuadruplicaba el nivel provincial, mientras hoy lo duplica. Lo que ocurrió en doce años no fue que La Caldera se llenara de casas vacías: fue que multiplicó por 1,9 su parque habitacional \textit{manteniendo} una proporción de viviendas sin hogar que ya era alta, mientras el resto de la provincia se le acercaba.

Esto obliga a reformular el argumento del capítulo, no a abandonarlo. La hipótesis del propietario no residente como novedad del ciclo inmobiliario reciente queda descartada. Lo que el dato sostiene es algo distinto y de más largo aliento: el departamento tiene, desde al menos 2010, una estructura de tenencia en la que una de cada cinco viviendas no aloja a nadie de manera permanente, y el ciclo de loteos reprodujo esa estructura a escala doble en lugar de corregirla.

Dicho de otro modo: el corredor norte no se convirtió en zona de segunda residencia; ya lo era, y lo que hizo el mercado fue duplicar su tamaño sin cambiar su composición.

\begin{cautela}
Corresponde una advertencia técnica sobre estas cifras, porque la comparación intercensal no es limpia y el libro no debe presentarla como si lo fuera.

Los datos de hogares de 2010 provienen de la base del \textbf{cuestionario ampliado}, que releva únicamente viviendas ocupadas y por lo tanto no permite contar directamente las deshabitadas; el total de viviendas particulares de 2010 proviene, en cambio, del recuento completo. Se están combinando dos fuentes, y ese cruce es válido sólo si el ampliado cubrió la totalidad de los hogares del departamento. Hay un indicio fuerte de que así fue: los 2.088 hogares, multiplicados por su distribución de tamaño, reconstruyen una población de 7.763 personas, exactamente la población departamental de 2010.

Existe además una medición alternativa. Si en lugar de hogares se usan las viviendas efectivamente ocupadas del ampliado --- 1.846 en La Caldera y 268.655 en la provincia ---, las tasas de 2010 resultan del 28,39 por ciento para el departamento y del 14,76 para la provincia. La razón entre ambas cae a 1,9. La diferencia entre las dos mediciones se explica porque en 2010 era frecuente que una vivienda alojara más de un hogar.

Bajo cualquiera de las dos mediciones se sostiene lo mismo: en 2010 el departamento ya estaba muy por encima del nivel provincial. Lo que varía es cuánto. La cifra definitiva requiere procesar la base del \textbf{cuestionario básico} de 2010 con la variable de condición de ocupación, que es un recuento completo y no exige combinar fuentes.

Un procesamiento adicional confirma la lectura de esa base. El promedio de personas por vivienda que el cuestionario ampliado arroja para el departamento en 2010 es de \textbf{4,19}, cifra que se aproxima a la de dividir la población departamental completa --- 7.763 --- por las 1.846 viviendas del ampliado, que da 4,21. Es decir que el ampliado relevó a toda la población del departamento distribuida en viviendas ocupadas: su universo de viviendas es el de las ocupadas, y no el total. La advertencia queda verificada. Y hay una más, de método censal: el Censo 2010 fue \textit{de hecho} ---contó a las personas donde pasaron la noche--- y el de 2022 fue el primero \textit{de derecho} ---las contó donde residen habitualmente---. La comparación de tasas entre ambos cruza ese cambio.
\end{cautela}

Ese mismo procesamiento agrega un dato de contexto que conviene retener porque despeja una explicación alternativa.

En 2010, el promedio de personas por vivienda ocupada de La Caldera --- 4,19 --- era el \textbf{cuarto más bajo de los veintitrés departamentos de Salta}, sólo por encima de Guachipas (3,71), Rosario de la Frontera (4,09) y La Candelaria (4,15), y muy por debajo de Chicoana (5,10), Los Andes (5,05) o Cafayate (5,01). El promedio simple departamental era de 4,56 y el provincial ponderado, de aproximadamente 4,52.

En 2022 la posición relativa se mantiene: 3,06 personas por hogar en el departamento contra 3,43 en la provincia. En ambos censos las unidades caldereñas son más chicas que el promedio provincial: un 7 por ciento en 2010, medido por vivienda ocupada, y un 11 por ciento en 2022, medido por hogar.

Eso importa por lo siguiente. Un hogar más chico reduce la relación entre población y viviendas sin que haya una sola casa vacía, y podría por sí solo explicar parte de la brecha. Pero la magnitud no alcanza: la diferencia de tamaño del hogar es de alrededor del diez por ciento y es estable entre censos, mientras que la diferencia en la tasa de viviendas sin hogar es de dos a cuatro veces. El tamaño del hogar explica una fracción menor de la brecha; la vivienda sin residentes explica el resto.



\section{Serie histórica}

\begin{longtable}{@{}p{3.2cm}rrrr@{}}
\toprule
\textbf{Unidad} & \textbf{1991} & \textbf{2001} & \textbf{2010} & \textbf{2022} \\
\midrule
\endhead
Vaqueros (localidad) & 1.877 & 2.980 & --- & --- \\
La Caldera (localidad) & 806 & 1.565 & --- & --- \\
\addlinespace
Departamento La Caldera & --- & --- & 7.763 & \textbf{12.299} \\
Viviendas particulares, departamento & --- & --- & 2.578 & \textbf{4.991} \\
\addlinespace
Provincia de Salta & --- & --- & 1.214.441 & 1.441.351 \\
\bottomrule
\end{longtable}

La serie está incompleta y la tabla lo muestra en lugar de disimularlo. Las cifras de 1991 y 2001 corresponden a localidades y las de 2010 y 2022 al departamento: no son sumables ni comparables directamente, porque el departamento incluye población dispersa fuera de ambos ejidos urbanos.

Hay, en cambio, un punto de partida comparable, y mucho más antiguo. El IV Censo General de la Nación registró en 1947 \textbf{2.831 habitantes} en el departamento Caldera ---que entonces, como hoy, incluía a Vaqueros--- sobre 290.826 de la provincia (Tomo I, p.\ 322, cuadro 7). Entre 1947 y 2022 el departamento multiplicó su población por \textbf{4,3}; la provincia, por \textbf{5,0}. Lo que distingue al departamento es el crecimiento reciente, no el del largo plazo.

Lo que sí puede afirmarse, con las unidades comparables, es la magnitud del período reciente: el departamento sumó 4.536 habitantes y 2.413 viviendas particulares en doce años. Sobre un municipio de tercera categoría.

\section{La categoría se fijó el 9 de diciembre de 1931}

Este capítulo sostiene que la clasificación municipal del departamento está congelada. Ahora
puede fecharse el día en que se fijó.

\begin{ficha}{Decreto Nº 14.205 del 9 de diciembre de 1931}
Dictado por el \textbf{Interventor Nacional} Mariano Gómez en acuerdo de ministros,
invocando los artículos 171 y 204 de la Constitución provincial, la Ley Orgánica de
Municipalidades del 20 de diciembre de 1898 y la Ley de Elecciones Municipales de 1919.

Los considerandos dejan constancia de que «esta Intervención Nacional ha practicado las
investigaciones a que se refiere el Art.\ 10º de la Ley Orgánica de Municipalidades \ldots\
y en su mérito \textbf{ha clasificado la categoría de los diversos Municipios de la
Provincia} en la forma que a cada uno corresponde».

\textbf{Art.\ 1º:} fija la categoría de los municipios existentes. La Capital, primera
categoría. \textbf{Caldera, tercera.}

\textbf{Art.\ 3º:} las de primera y segunda categoría se componen de un Concejo Deliberante
de nueve y cinco concejales elegidos por el pueblo, y de un \textbf{Intendente nombrado y
removido por el Poder Ejecutivo con acuerdo del Senado}. Las de tercera, «por Comisiones
compuestas de tres a cinco miembros, según sea su población, elegidos directamente por el
pueblo, de los cuales uno, que desempeñará el cargo de \textbf{Presidente, será nombrado por
el Poder Ejecutivo}».

\textbf{Art.\ 4º:} los concejales y miembros de comisión duran dos años; los intendentes,
cuatro; \textbf{los presidentes de comisión, uno}.
\end{ficha}

Ese mandato de un año se cumplía. En marzo de 1933, «habiendo don \textbf{Augusto Regis\index{Regis, Augusto}}
terminado su período constitucional en las funciones de Presidente de la Comisión Municipal
de La Caldera», el Poder Ejecutivo lo volvió a nombrar en el mismo cargo, invocando el
artículo 182 de la Constitución provincial.

Es el mismo Augusto Regis\index{Regis, Augusto} que en 1908 figura como jurado suplente del departamento, que en
1910 integró la comisión de siete vecinos para las defensas contra el río y que en 1922 fue
designado por el Interventor Nacional miembro de esa comisión municipal. Veinticinco años
después seguía en el cuerpo, ahora presidiéndolo por designación anual del gobernador.

Tres consecuencias, y la primera reordena lo que este capítulo venía diciendo.

\textbf{La categorización no es un error administrativo ni un olvido: es un acto.} Tiene
fecha, número, firmante y considerandos que declaran haberse practicado las investigaciones
del artículo 10º de la ley orgánica. Alguien clasificó a La Caldera en tercera categoría el
9 de diciembre de 1931, con fundamento y por escrito.

\textbf{Y en 1931 la clasificación era correcta.} El departamento tenía entonces unos pocos
miles de habitantes y la tercera categoría le correspondía sin discusión. Lo que este libro
documenta no es que la clasificación haya nacido mal, sino que \textbf{nunca volvió a
hacerse}. El decreto que la fijó habla de «los Municipios existentes a la fecha de este
decreto»: era una foto, y quedó como retrato permanente.

\textbf{Y en 1947 la ley dijo expresamente que la foto seguiría.} La Ley 2131 (original 853), que ese año fijó los límites de todos los distritos municipales de la provincia y asignó a La Caldera «todo el territorio del Departamento», dispuso en su artículo 2º: «Hasta tanto se realice un censo general de la población las municipalidades mantendrán sus actuales categorías y organización» (sancionada el 24 de julio de 1947 y publicada en el Boletín Oficial Nº 2.912 del 7 de agosto). El congelamiento, entonces, no fue sólo inercia: quedó escrito como regla provisoria, atado a un censo que la provincia hizo después varias veces sin que la categoría se revisara.

\textbf{Y la tercera consecuencia es sobre el cargo.} El artículo 3º explica por qué el
titular del ejecutivo municipal se llamaba presidente y no intendente, y por qué lo nombraba
el gobernador: no era una excepción de la dictadura ni un rasgo de la intervención de 1922,
era el régimen ordinario de los municipios de tercera categoría. Y duraba \textbf{un año},
contra los cuatro del intendente de primera o segunda. En 1980, la dictadura llevó ese régimen a todos los municipios: la Ley de facto 5686 hizo designar y remover por el Ejecutivo a todos los intendentes y presidentes de comisión, les dio también las funciones de los concejos, con autorización previa del gobierno para las decisiones de fondo, y creó «Consejos Asesores Comunales» de vecinos nombrados por el ministerio, con un mínimo de tres miembros en los municipios de tercera categoría.

De modo que la asimetría que este libro describe entre La Caldera y la capital no empieza en
el mercado de suelo ni en el crecimiento demográfico. \textbf{Está inscripta en el diseño
institucional desde 1931}, y la clasificación que ese instrumento inscribió es la que la estadística
provincial todavía usa noventa y cinco años después, aunque la ley que le daba sustento se derogó en 2018.

\section{No es sólo el municipio: es toda la clasificación}

En 1934, en el término de cinco días, el Poder Ejecutivo provincial clasificó dos servicios
más del departamento, y en los dos lo puso en el escalón más bajo.

El \textbf{Decreto Nº 17.606}, del 12 de marzo, fija la categoría de las oficinas del
Registro Civil de la provincia según el movimiento estadístico-demográfico de 1933. La
Caldera figura en el orden 33 y queda en \textbf{tercera categoría, la última de tres}.

El \textbf{Decreto Nº 17.640}, del 17 de marzo, fija la de las Receptorías de Rentas de
campaña según lo recaudado por cada una durante 1933. La de La Caldera queda en
\textbf{quinta categoría, la última de cinco}.

Sumadas a la clasificación municipal de 1931 --- tercera de tres ---, son tres actos
independientes, de tres ministerios distintos, con tres criterios distintos --- población,
demografía registrada y recaudación ---, y los tres colocan al departamento en el último
escalón disponible.

\begin{cautela}
Conviene no leer esto como una decisión concertada, porque no lo fue. Cada uno de esos
decretos clasifica a todos los distritos de la provincia según un criterio explícito, y La
Caldera queda abajo porque en 1934 era efectivamente un departamento chico, con poca
población registrada y poca recaudación. Las tres clasificaciones eran correctas.

Lo que interesa es el \textbf{efecto acumulado y su permanencia}. Cada clasificación
determina dotación de personal, presupuesto y jerarquía del funcionario a cargo: una oficina
de tercera tiene menos gente y menos rango que una de primera, y una receptoría de quinta,
menos aún. De modo que el departamento no recibió una asignación baja sino \textbf{tres,
simultáneas y por vías separadas}, y cada una reforzaba a las otras: menos recaudación
justificaba menos receptoría, y menos oficina producía menos registro.

Es el mismo mecanismo que este capítulo describe para la categoría municipal, multiplicado.
Y con la misma característica: se fijó una vez, sobre una foto correcta, y quedó.
\end{cautela}

\section{Y noventa y un años después sigue en tercera}

El informe de gobiernos locales de la Dirección General de Estadísticas provincial clasifica a los municipios salteños en tres categorías. \textbf{La Caldera y Vaqueros son ambos M3: municipios de tercera categoría}, la más baja, junto con localidades como Tolar Grande, La Poma, Iruya o Nazareno.

\begin{cautela}
El departamento que más creció de la provincia y el cuarto del país figura, en la estadística oficial, administrado por dos municipios de la categoría más baja.

Durante casi nueve décadas esa categoría determinó forma de gobierno, atribuciones y duración de los cargos. Hoy no determina nada, y eso es lo que este apartado tiene que explicar: la ley que la sostenía ya no existe.
\end{cautela}

\subsection*{La ley que fijaba las categorías, y la que la reemplazó}

Hasta 2018 los criterios los fijaba la Ley provincial 1349, Ley Orgánica de Municipalidades.

\begin{ficha}{Categorías municipales, Ley 1349 (derogada en 2018)}
\begin{itemize}[leftmargin=*,nosep]
\item \textbf{Primera categoría} (art.\ 4): más de diez mil habitantes. Gobierno de Intendente y Concejo Deliberante de \textbf{nueve} concejales (art.\ 5).
\item \textbf{Segunda categoría} (art.\ 4): población que exceda de cinco mil habitantes. Intendente y Concejo de \textbf{cinco} concejales (art.\ 5).
\item \textbf{Tercera categoría} (art.\ 32): entre quinientos y cinco mil habitantes, distribuidos en un radio no mayor de tres kilómetros cuadrados, sin otro centro urbano organizado como municipio en quince kilómetros. Se denominaban \textbf{Comisiones Municipales}, de tres a cinco miembros, y \textbf{su presidente era nombrado y removido por el Poder Ejecutivo} (art.\ 33).
\end{itemize}
El artículo 2 establecía que la categoría se determinaba \textbf{por ley}, y el artículo 3 que, comprobado por un censo que un municipio había alcanzado el número de habitantes para ascender, \textbf{el Poder Ejecutivo enviaría a la Legislatura el proyecto de ley correspondiente}.
\end{ficha}

El 13 de noviembre de 2018 la Legislatura sancionó la \textbf{Ley 8126 de Régimen de Municipalidades} (B.O.\ 17/12/2018), que rige para los municipios que no dictaron carta orgánica y cuyo artículo 93 deroga la 1349. \textbf{La ley nueva no tiene categorías.} Todo municipio al que se aplica se gobierna por un Concejo Deliberante y un intendente elegidos por el pueblo (art.\ 45), con las mismas atribuciones (arts.\ 60 y 69). La única distinción que subsiste es la que separa a los municipios que dictaron su carta orgánica de los que no. El Gobierno provincial presentó entonces la reforma como el reemplazo de una norma de los años treinta que clasificaba a los municipios en categorías y los dejaba sin autonomía, y precisó que regiría para los cuarenta municipios sin carta orgánica, los que no superaban los diez mil habitantes.

\begin{cautela}
Esto resuelve la discrepancia que una versión anterior de este capítulo dejaba abierta, y la resuelve en una dirección que no se esperaba.

El informe de gobiernos locales \textbf{no usa M3 como abreviatura de tamaño}: lo define expresamente como ``municipio de 3\textordmasculine{} categoría'', y ubica en esa lista a La Caldera y a Vaqueros junto con otros veintisiete, entre ellos San Lorenzo y Coronel Moldes. La tercera categoría era la Comisión Municipal de la Ley 1349, con presidente designado por el Poder Ejecutivo; desde 2018 ni esa figura ni ninguna categoría existen en el régimen provincial. \textbf{La clasificación estadística no corresponde a ninguna ley vigente.} Es la huella de una derogada.

Y no es la huella de su versión final, sino de su foto más vieja. De los veintinueve municipios clasificados M3, \textbf{seis superan los cinco mil habitantes} que la 1349 fijaba como techo de la tercera categoría ---cifra que no nace con esa ley: el capítulo \ref{cap:siglo} documenta que ya en enero de 1919 un decreto provincial la atribuía al artículo 175 de la Constitución como condición de las municipalidades autónomas---:
\end{cautela}

{\small
\begin{longtable}{@{}lr@{}}
\toprule
\textbf{Municipio clasificado M3} & \textbf{Población 2022} \\
\midrule
\endhead
\textbf{San Lorenzo} & \textbf{32.932} \\
Vaqueros & 7.774 \\
General Pizarro & 7.244 \\
San Antonio de los Cobres & 6.853 \\
Cachi & 6.791 \\
Coronel Moldes & 5.695 \\
\bottomrule
\end{longtable}}

\begin{cautela}
San Lorenzo es el caso límite: con \textbf{32.932 habitantes} triplica el umbral de la primera categoría de la ley derogada, dictó su propia carta orgánica en 2018, según el mismo parte oficial ---de modo que ni siquiera la Ley 8126 lo alcanza--- y figura igual como M3.

El capítulo \ref{cap:politica} fecha el origen. En 1983 la clasificación era \textbf{exacta}: el Boletín Oficial de ese año registra a La Caldera y a Vaqueros como \textbf{Comisiones Municipales}, con presidentes designados cuyas renuncias el Poder Ejecutivo aceptó el 9 de diciembre de 1983, en decretos distintos de los que aceptaban las renuncias de los intendentes de las Municipalidades. La categoría que hoy comparten con San Lorenzo es una clasificación de 1931, que la ley de 1947 congeló hasta un censo, que en 1983 todavía era exacta, que nadie actualizó mientras la ley la exigía, y que \textbf{sobrevive a la derogación de esa ley}.

Para este libro la consecuencia es doble. La primera es de método: la categoría \textbf{no mide} capacidad institucional ni forma de gobierno ---la ley vigente iguala a todos los municipios sin carta orgánica---, de modo que ninguna comparación de este libro se apoya en ella. La segunda es sustantiva, y más incómoda: lo que la categoría sí mide, con exactitud involuntaria, es la antigüedad del dato con que la provincia describe a sus gobiernos locales. \textbf{La estadística oficial de 2022 clasifica al departamento que más creció de la provincia con una categoría de 1931 y una ley que dejó de regir en 2018.}
\end{cautela}

Y un dato de método que el mismo informe entrega al pasar, con consecuencias que exceden a este capítulo.

La geometría de los gobiernos locales que el censo utiliza fue construida por la Dirección General de Inmuebles \textbf{de acuerdo con la Ley provincial Nº 2131 de 1947}. El informe advierte, en notas al pie, que siete municipios --- entre ellos La Merced, General Güemes, Cerrillos, Campo Santo y El Bordo --- ``incluyen población y viviendas de un área cuyos límites difieren de los establecidos'' por esa ley.

\textbf{La Caldera y Vaqueros no están entre esos siete.} La ley de 1947 asignaba todo el departamento a La Caldera, de modo que la separación entre ambos municipios sólo puede provenir de las leyes posteriores ---la 4365 de 1970 y la 4987 de 1974---, que el informe no menciona. Qué versión del límite se usó para separar sus cifras es algo que el informe no aclara. Toda comparación intermunicipal de este libro descansa sobre esa geometría, y conviene decirlo antes de comparar.


\section{Un enclave de fuerza de trabajo}

El Censo 2010 aporta un último elemento que ayuda a caracterizar a la población que ocupa esas viviendas.

\begin{ficha}{Condición de actividad, Censo 2010, población de 14 años y más}
\begin{itemize}[leftmargin=*,nosep]
\item \textbf{Departamento La Caldera:} 64,37\% de ocupados, 3,24\% de desocupados y 32,39\% de inactivos, sobre 5.555 personas.
\item \textbf{Provincia de Salta:} 54,42\% de ocupados, 4,49\% de desocupados y 41,09\% de inactivos.
\item La Caldera registra \textbf{la tasa de ocupación más alta de los veintitrés departamentos de Salta} y \textbf{la tasa de inactividad más baja}. La siguen Cafayate (60,21\%), Cachi (58,89\%), Cerrillos (58,53\%) y Capital (58,20\%). En el extremo opuesto: Santa Victoria (28,95\%), Iruya (38,27\%) y Molinos (40,92\%).
\item Medido sobre población total, el departamento tenía 46,06\% de ocupados contra 38,22\% provincial.
\end{itemize}
\textit{Fuente: procesamiento Redatam, Censo Nacional 2010, cuestionario ampliado con expansor.}
\end{ficha}

La estructura de edades no explica la diferencia: la población de catorce años y más representa el 71,6 por ciento del departamento y el 70,2 de la provincia, prácticamente lo mismo. Lo que difiere es la relación de esa población con el trabajo.

Un departamento con la ocupación más alta y la inactividad más baja de la provincia, con hogares alrededor de un diez por ciento más chicos que el promedio y con una de cada cinco viviendas sin residentes permanentes, describe un perfil poblacional preciso: adultos en edad activa, empleados, en hogares reducidos, sobre un parque habitacional que excede a la población residente.

Es, en datos censales, la definición de ciudad dormitorio. La población del departamento trabaja --- y en su mayoría, dado que la base productiva local es escasa, trabaja en la capital, a una docena de kilómetros de Vaqueros y unos veinte de La Caldera, por la única vía de acceso que el plan urbano de 2015 señalaba como principal problema del ejido.

Conviene poner ese perfil junto al artículo 14 del Código de Vaqueros, que el capítulo \ref{cap:vaqueros} documentó. Esa norma establece que la condición periférica metropolitana del municipio ``no implica su consideración como mera ciudad dormitorio de Salta capital''. El Censo de 2010, anterior a cualquiera de las versiones de esa norma, ya describía exactamente eso.

Corresponde consignar además, porque afecta la lectura de dos procesamientos anteriores, que la columna ``Total seleccionado'' que aparece sin universo declarado en las salidas de Redatam del Censo 2010 corresponde a la \textbf{población ocupada}. Este libro había señalado que esa columna era ininterpretable y se había abstenido de usarla; la presente consulta la identifica y permite recuperarla.
\section{Agosto de 2024: los propietarios que no viven allí}

En agosto de 2024, la Dirección de Administración del municipio emitió un recordatorio dirigido a \textbf{los propietarios de inmuebles registrados en el municipio pero con domicilio fiscal en otras localidades}, para que regularizaran sus deudas tributarias. Según \textit{Nuevo Diario de Salta}, la medida respondió a un incremento de la morosidad y comprendía tasa de servicios públicos, impuesto inmobiliario, impuesto automotor y actividades comerciales. Se anunció el inicio de intimaciones formales con sanciones e intereses punitorios.

Esta pieza cierra el argumento del artículo 30 del Código.

El Código manda gravar de manera permanente los lotes no edificados y los loteos en áreas no aptas. El intendente sostuvo en enero de 2026 que las propiedades de los barrios inundados no pagan impuestos municipales. Y en agosto de 2024 el propio municipio reconoció públicamente que existe un universo identificable de propietarios con inmueble en La Caldera y domicilio fiscal fuera, en morosidad creciente.

Ese universo tiene un nombre en la descripción de este libro: es el comprador del loteo periurbano. Adquiere una parcela en el corredor norte, no reside allí, no edifica o edifica para uso de fin de semana, y tributa poco o nada. El Censo permite ahora dimensionar el fenómeno: hay 988 viviendas particulares sin hogar residente en el departamento, casi una de cada cinco. Lo que el Censo todavía no permite es distinguir, dentro de esas 988, cuántas son segundas residencias, cuántas están en venta y cuántas en construcción. El municipio le provee camino, alumbrado y recolección --- los servicios mínimos que el artículo 42 asigna a los núcleos no planificados --- y cuando el río crece, le provee además maquinaria, defensas y asistencia.

De modo que la fórmula puede enunciarse con precisión. La conversión de monte en parcela produce, para el desarrollador, una renta que se realiza y se retira; para el comprador, un patrimonio que se valoriza sin cargas proporcionales; y para el municipio, una obligación de servicio permanente sobre suelo que no genera recaudación equivalente.

No es que el municipio no cobre porque no quiere: su base imponible está compuesta, en proporción creciente, por contribuyentes que no viven allí. Y su capacidad técnica está determinada por esa base.


\section{Quién vive en cada radio: la formación}

El Censo 2022 permite bajar por debajo de la fracción, hasta el \textbf{radio censal}, una unidad de pocas manzanas en las zonas urbanas. El departamento tiene veintiocho: quince en la fracción 01, que es Vaqueros; nueve en la 02, que es el pueblo de La Caldera y su entorno; y cuatro en la 03, al este del municipio. Suman 12.246 personas en viviendas particulares, la misma cifra que usa este capítulo.

La variable disponible es el \textbf{nivel más alto que cursó} cada persona. Con ella y con la estructura de edades del mismo radio se construye un indicador sencillo: cuántas personas cursaron un nivel universitario o de posgrado por cada cien habitantes de veinticinco años o más.

\begin{ficha}{Formación por área --- Censo 2022}
\begin{itemize}[leftmargin=*,nosep]
\item \textbf{Provincia de Salta:} 10,8 personas con universitario o posgrado cursado por cada cien adultos; 23,9 si se suma el terciario.
\item \textbf{Departamento Capital:} 16,7 y 30,9.
\item \textbf{Departamento La Caldera:} \textbf{25,5 y 40,6}.
\item \textbf{Fracción 01 (Vaqueros):} 30,3 y 44,3. \textbf{Fracción 02 (pueblo de La Caldera):} 13,6 y 30,4. \textbf{Fracción 03:} 33,0 y 49,4.
\end{itemize}
\end{ficha}

\textbf{El departamento tiene, en proporción, más población con estudios universitarios que el departamento Capital}: una vez y media la de la Capital y más del doble de la provincial. Es el mismo departamento que el capítulo \ref{cap:politica} muestra con un padrón electoral mayor que su población, y que este capítulo describe con una de cada cinco viviendas sin residentes permanentes. Los tres datos dibujan la misma figura: \textbf{una extensión residencial de la capital}, con su perfil educativo.

Pero el promedio esconde dos departamentos. La lámina \ref{fig:formacionradios} lo muestra radio por radio.

\begin{figure}[p]
\centering
\includegraphics[width=\textwidth,height=0.86\textheight,keepaspectratio]{fig-formacion-radios.png}
\caption[Formación universitaria por radio censal, 2022]{\textbf{Formación universitaria por radio censal, departamento La Caldera, 2022.} A: el departamento; los recuadros B y C marcan los detalles. B: el pueblo de La Caldera y la matrícula 4366, con la línea de ribera del pueblo (Resolución 023/14) y las del río La Caldera y el arroyo El Manzano sobre esa finca (Resolución 185/14), trazadas con sus coordenadas publicadas; el triángulo es el primer punto de la margen derecha del río. C: el área urbana de Vaqueros. En gris, los radios de otros departamentos. Elaboración propia con datos del INDEC (Censo 2022, Redatam) y la cartografía de radios corregida por G.~M.~Rodríguez (CEUR-CONICET, versión 1.0); coordenadas de los edictos de las Resoluciones 023/14 y 185/14.}\label{fig:formacionradios}
\end{figure}

\begin{cautela}
\textbf{Y hay un cruce que este libro tenía en tres capítulos distintos y que conviene hacer acá, porque cambia lo que la fracción 03 significa.}

\begin{center}\small
\begin{tabular}{@{}p{7.6cm}r@{}}
\toprule
Viviendas sin hogar residente & \textbf{43,17\,\%} \textit{(5º de 159 fracciones)} \\
Personas por vivienda & \textbf{1,39} \\
Formación universitaria o de posgrado & \textbf{33,0} por cada cien adultos \\
Hogares sin agua de red \textit{(cap.\ \ref{cap:aguabaja})} & \textbf{58,0\,\%} \\
Población de 15 a 29 años \textit{(cap.\ \ref{cap:ausencias})} & \textbf{14,4\,\%} \\
\bottomrule
\end{tabular}
\end{center}

\textbf{Es el nivel de formación más alto del departamento} ---por encima de Vaqueros, que marca 30,3, y muy por encima del pueblo de La Caldera, que marca 13,6--- \textbf{sobre la peor cobertura de servicios y la mayor proporción de casas vacías.}

Setecientas sesenta y cuatro personas, cuatro de cada diez viviendas sin nadie adentro, 1,4 habitantes por vivienda, seis de cada diez hogares sin agua de red y un tercio de los adultos con paso por la universidad \textbf{no describen un sector rural despoblado}. Describen la segunda residencia de sectores medios y altos que este capítulo tiene por objeto, y que el remate de la casa quinta de Lesser documenta desde 1927.

El capítulo \ref{cap:ausencias} llegó a la misma conclusión por descarte ---«un sector rural de baja densidad con vivienda de fin de semana»--- sin este dato. Con él, la conclusión deja de ser una inferencia sobre lo que falta y pasa a apoyarse en lo que sobra. \textbf{Y ubica el fenómeno}: la vacancia del departamento no está repartida ni concentrada en el pueblo; está en el extremo oriental, sobre la Ruta Provincial 11, que es donde el apéndice \ref{cap:pedidos} ubica los dos fraccionamientos de El Gallinato que el capítulo \ref{cap:opacidad} observa en el visor y que siguen sin identificar.

\textbf{Con la reserva de tamaño de siempre}: son 764 personas y 549 viviendas, y en un universo así un puñado de casos mueve varios puntos. La dirección del hallazgo no depende de eso; su magnitud, sí.
\end{cautela}

\textbf{En Vaqueros}, nueve de los quince radios superan los treinta universitarios por cada cien adultos, y el 0104 llega a cincuenta. Sólo dos quedan por debajo del promedio provincial: el 0113 y el 0101, un radio rural de setenta y ocho habitantes.

\textbf{En el pueblo de La Caldera} la distribución es otra. Los radios del casco oscilan entre 13 y 19, y el más poblado del municipio, el 0207, con 1.279 habitantes, tiene \textbf{5}: menos de la mitad del promedio provincial. Al sur del pueblo, en cambio, los radios 0208 y 0209 marcan 21 y 24, por encima de la capital. \textbf{En el 0209, más de la mitad de los adultos cursó un terciario o una carrera universitaria.}

\textbf{El Durazno está en el 0208 o en el 0209, y la coordenada no permite decir en cuál.} La línea de ribera del río La Caldera sobre la matrícula 4366 ---la finca que el edicto de agua de 2015 vincula con el loteo, aprobado sobre la 6.558--- recorre los dos radios: de los catorce puntos de su margen derecha, publicados en el edicto de la Resolución 185/14, once caen en el 0208 y los tres últimos, río abajo, en el 0209; la del arroyo El Manzano, sobre la misma finca, cruza el 0203 y el 0208. Una versión anterior de este párrafo ubicaba el primer punto en el 0209; trazado con las coordenadas del edicto sobre la cartografía censal, cae en el 0208. El 0208, entre el 0209 y el pueblo, coincide con la ubicación que el Código de Ordenamiento da a El Nogalar y a Potreros de La Caldera, al sudoeste del casco. No es un radio envejecido: el 28 por ciento de su población tiene menos de quince años y sólo el 6 por ciento más de sesenta y cinco.

La lámina \ref{fig:formacioncomparacion} reúne las comparaciones.

\begin{figure}[htbp]
\centering
\includegraphics[width=\textwidth]{fig-formacion-comparacion.png}
\caption[Formación por área y radios seleccionados, 2022]{\textbf{Formación por área y radios seleccionados, 2022.} Personas que cursaron cada nivel por cada cien habitantes de veinticinco años o más. Debajo de la línea punteada, las áreas; encima, cuatro radios del departamento. Elaboración propia con datos del INDEC (Censo 2022, Redatam).}\label{fig:formacioncomparacion}
\end{figure}

\begin{cautela}
Cuatro límites, y ninguno es menor.

\textbf{Cursar no es terminar.} La variable registra el nivel más alto al que la persona asistió, no si lo completó. Quien dejó la universidad en primer año cuenta igual que quien se recibió.

\textbf{El indicador es aproximado.} Edad y nivel vienen en el mismo cuadro, pero no cruzados: el numerador incluye a los estudiantes universitarios de dieciocho a veinticuatro años, y el denominador sólo a los mayores de veinticinco. Sirve para comparar áreas entre sí, no como tasa exacta.

\textbf{Los radios son chicos.} El 0209 tiene 150 habitantes y el 0208, 125. Pocas familias mueven mucho el porcentaje. Y en esos loteos buena parte de las viviendas son de fin de semana: sus dueños no se cuentan allí.

\textbf{La identificación de los barrios es una inferencia.} Las coordenadas publicadas muestran que la finca 4366 abarca el 0208 y el 0209, pero ninguna señala el loteo; la ubicación de los barrios descansa en la posición relativa que les da el Código de Ordenamiento. Cuál de los dos contiene a El Nogalar y cuál a Potreros de La Caldera no lo establece este libro.
\end{cautela}

\section{Vacancia y turismo son el mismo uso del suelo}

El 18 de febrero de 2026, el mismo mes en que se aprobó el Plan de Manejo, el municipio de La Caldera publicó un parte sobre el fin de semana de Carnaval.

\begin{ficha}{Parte municipal, 18 de febrero de 2026 --- área Turismo}
Tasa Neta de Ocupación Turística promedio del \textbf{80,9 por ciento}, según datos del Ministerio de Turismo y Deportes de la Provincia, que ubican a La Caldera ``entre los principales municipios de la provincia con más alta ocupación''. En los registros locales se contabilizaron \textbf{cuatro establecimientos con el 100 por ciento} de ocupación y otros con el 80 por ciento.

El intendente destacó que hubo ``alta ocupación de \textbf{cabañas}, algunas al 100 por ciento'', y que el municipio obtuvo dos premios turísticos: uno provincial, con el que ``está creando un nuevo circuito turístico en la zona del dique'', y otro nacional para equipamiento tecnológico.
\end{ficha}

\begin{cautela}
Léanse los dos datos juntos, porque son del mismo municipio, del mismo mes y de la misma superficie.

El municipio de La Caldera tiene la \textbf{tercera menor ocupación por vivienda de la provincia} --- 2,23 personas por vivienda particular contra 3,13 provinciales --- y más de una de cada cuatro viviendas sin residentes permanentes (27,9 por ciento en sus fracciones 02 y 03; una de cada cinco es la cifra del departamento). Y en febrero de 2026 informa, con estadística provincial, una ocupación turística del 80,9 por ciento con establecimientos al tope.

No hay contradicción entre ambas cifras: son la misma cosa vista desde dos oficinas. Un parque construido que está vacío la mayor parte del año y lleno los fines de semana largos no es un parque residencial subutilizado; es un parque de uso intermitente. La estadística de vivienda lo registra como vacancia y la estadística de turismo lo registra como éxito.

\textbf{Con una advertencia que corresponde hacer.} El dato turístico mide plazas de alojamiento comercial --- cabañas, hosterías, establecimientos registrados ---, no viviendas particulares vacías. Son universos distintos y este libro no los confunde: la ocupación del 80,9 por ciento no dice nada directo sobre las viviendas que el censo cuenta como deshabitadas. Lo que en cambio muestra es que existe, documentada por el propio Estado, \textbf{una demanda de uso temporario de suelo residencial en el departamento}, y que el municipio la reconoce como su activo económico y la promueve --- incluido un nuevo circuito en la zona del dique, que es el área que la Ley 4486 expropió en 1972.

Nada de esto es un reproche. Un municipio de cuatro mil quinientos habitantes que encuentra en el turismo una fuente de ingreso hace lo que puede hacer. Lo que este libro registra es que \textbf{la vacancia y el turismo no son dos problemas sino dos nombres del mismo uso del suelo}, y que la política municipal fomenta uno mientras la estadística censal mide el otro, sin que ninguna oficina tenga que mirar las dos cifras a la vez.
\end{cautela}


\subsection*{Y el Estado subsidia la capacidad turística}

El fomento del uso temporario no es sólo discurso municipal. Tiene instrumento provincial, y en el departamento se aplicó.

\begin{ficha}{Decreto Nº 3069 del 14 de octubre de 2014 --- B.O.\ Nº 19.409}
Expediente 136-135.448/2013. Ratifica el Contrato de Promoción Turística suscrito el 24 de septiembre de 2014 entre la Provincia y la firma \textbf{La Caldera S.A.}, en el marco de la \textbf{Ley 6064} de Promoción de la Actividad Turística y su modificatoria 7281.

\textbf{Proyecto:} ``Ampliación Hostería La Caldera'' --- construcción, equipamiento y explotación de \textbf{dos nuevas alas de dormitorios y un Salón de Usos Múltiples}, en inmueble propio de 12.427,13 m², \textbf{matrícula catastral 3.196}, parcela 5, manzana 1, sección D, localidad de La Caldera.

\textbf{Beneficios que el régimen prevé:} exención de todos o algunos de los tributos provinciales vigentes o a crearse, con exclusión de las tasas retributivas de servicios; y entrega de \textbf{certificados de crédito fiscal} aplicables al pago de los impuestos a las Actividades Económicas, a los Sellos y al Inmobiliario Rural.

Firman: Urtubey, Saravia, Parodi, Simón Padrós. \textbf{Los alcances concretos de los beneficios están en el Anexo}, que forma parte del decreto y no se publica en el cuerpo.
\end{ficha}

\begin{cautela}
Tres observaciones.

\textbf{La primera cierra el argumento de este apartado.} La ocupación turística que el municipio celebra en 2026 no creció sola: el Estado provincial la promovió con exenciones impositivas y créditos fiscales. En 2014, mientras los vecinos de la zona alta de Vaqueros esperaban la red de agua y el Ministerio exigía a un club de campo que previera abastecerse en forma privada, la Provincia eximía de tributos a una hostería para que ampliara sus dormitorios.

Ninguna de esas tres decisiones es objetable por separado, y las tres son políticas públicas legítimas. Juntas describen una asignación: \textbf{el Estado subsidió la capacidad de alojamiento temporario del mismo territorio cuya capacidad de alojamiento permanente no lograba dotar de agua}.

\textbf{La segunda es de método.} El decreto ratifica un contrato ``por el cual se fijan los alcances de los beneficios concedidos'', y ese contrato es el Anexo. \textbf{El Anexo no se publica.} Cuánto se eximió, por cuántos años y sobre qué base imponible no consta en el Boletín Oficial: consta en el expediente 136-135.448/2013. Es el mismo patrón que el capítulo \ref{cap:opacidad} describe --- el acto se publica, su contenido operativo no ---, y aquí recae sobre dinero público resignado.

\textbf{La tercera es un nombre.} La beneficiaria se llama \textbf{La Caldera S.A.} y es titular de un inmueble de 12.427 metros cuadrados en el pueblo, matrícula 3.196. Este libro no sabe nada más de ella: ni quiénes la integran, ni desde cuándo, ni qué otros inmuebles posee. Una sociedad anónima que lleva el nombre del municipio donde opera merece, por lo menos, que se busque su legajo.
\end{cautela}

\begin{cautela}
Y no fue un caso aislado. Cuatro meses después, el \textbf{Decreto Nº 463 del 9 de febrero de 2015} ratificó un segundo contrato de promoción turística, suscrito el 29 de diciembre de 2014 entre la Provincia y la firma \textbf{Proyecto 1 S.A.}, para el proyecto \textbf{``Ampliación Hotel La Caldera''}, en un inmueble propio de \textbf{matrícula 3.337}, con una superficie cubierta total de 1.605 metros cuadrados. Expediente 136-34167/2010. Se publicó en el Boletín Oficial Nº 19.485 del 19 de febrero de 2015 y, en lugar del contrato, lleva al pie la leyenda «anexo para su consulta en oficinas de nuestra repartición».

El decreto precisa el encuadre: la inversión entra en el artículo 3º de la Ley 6064 ``ya que corresponde a una \textbf{ampliación de más del 40\% de las existencias}''. Es decir, el beneficio impositivo se otorga a quien amplía en esa proporción lo que ya tiene.

Y la misma firma recibiría, seis meses más tarde, algo más que una exención: por el Decreto 2949 del 24 de agosto de 2015 la Provincia le entregó en \textbf{comodato gratuito por veinte años cuarenta y una hectáreas de tierra fiscal} para construir una cancha de golf de dieciocho hoyos, según se documenta en el capítulo \ref{cap:tierrafiscal}.

\textbf{Dos contratos de promoción turística en cuatro meses, sobre dos inmuebles distintos del mismo pueblo, con dos sociedades distintas.} Eso deja de ser un caso y pasa a ser una política: entre fines de 2014 y comienzos de 2015 el Estado provincial resignó tributos para expandir la capacidad hotelera de La Caldera.

Este libro no la objeta. Fomentar el turismo en un pueblo de cuatro mil quinientos habitantes es una decisión legítima y probablemente razonable.

El reglamento de esa ley ---Decreto 2448/85, B.O.\ Nº 12.356 del 4/12/1985--- permite medir cuánto control previó. Exige para cada solicitud la nómina de socios, la declaración jurada de bienes de los integrantes y el plan de inversiones (art.\ 24); obliga a los beneficiarios a presentar declaraciones juradas anuales bajo pena de caducidad de la exención (art.\ 32); manda a la Dirección Provincial de Turismo verificar el cumplimiento e informar al Ministerio \textbf{cada seis meses} (art.\ 33); y hace perder los beneficios si los bienes promovidos se transfieren, venden o arriendan, o si cambia la razón social sin autorización (art.\ 18). Admite además que los municipios adhieran eximiendo sus propias tasas y derechos de construcción (art.\ 37). De modo que la pregunta que este apartado deja sobre La Caldera S.A. ---quiénes la integran--- tiene respuesta escrita en su expediente de promoción, y la de qué se hizo con lo eximido, en los informes semestrales de Turismo. Si La Caldera adhirió al régimen, además, el pueblo resignó sus propias tasas sobre las mismas obras. Lo que registra es la \textbf{simultaneidad}, que ya no admite ser leída como coincidencia: en el mismo bienio, el Estado eximió de impuestos a dos hoteles, exigió a un club de campo que resolviera el agua por su cuenta, tramitó para otro un derecho permanente sobre un río, y dejó a la zona alta de Vaqueros con camión cisterna.
\end{cautela}


\subsection*{El sistema tarifario ya le puso precio a la vacancia}

El cuadro tarifario de la Compañía Salteña de Agua y Saneamiento aprobado por Resolución 2066/25 del Ente Regulador de los Servicios Públicos, con vigencia desde enero de 2026, contiene una categoría que este libro debe registrar.

\begin{ficha}{Categoría ``Baldío'' --- Cuadro tarifario COSAYSA, enero de 2026}
El cuadro distingue, además de las categorías residencial, comercial e industrial, una categoría \textbf{Baldío}, con cargo fijo propio por diámetro de conexión y precio propio por metro cúbico.

\textbf{Precio del metro cúbico, sólo agua:}
\begin{center}\small
\begin{tabular}{@{}lrr@{}}
\toprule
& \textbf{Residencial} & \textbf{Baldío} \\
\midrule
Zona 1 & \$342,53 & \textbf{\$404,01} \\
Zona 2 & \$258,57 & \textbf{\$330,39} \\
Zona 3 & \$191,58 & \textbf{\$293,76} \\
\bottomrule
\end{tabular}
\end{center}

En las tres zonas el metro cúbico del baldío \textbf{cuesta más} que el residencial: un 18, un 28 y un 53 por ciento más, respectivamente. En cargos fijos, el baldío de trece milímetros supera al residencial de categorías 2 y 3.
\end{ficha}

Hay tres cuestiones que se siguen de este cuadro, y conviene separarlas por grado de certeza.

\textbf{La primera es segura: existe una política tarifaria sobre el suelo no ocupado.} El sistema regulatorio provincial no trata al lote vacío como una ausencia; lo trata como una categoría de usuario, con tarifa propia y más cara que la del vecino que vive allí. Alguien decidió, y el Ente Regulador aprobó, que la conexión que no se usa pague más que la que se usa.

\textbf{La segunda es casi segura: hay un registro.} Para facturar una categoría hay que poder identificarla. La empresa sabe, conexión por conexión, cuáles son baldíos. Ese padrón es la contracara operativa del dato censal que este libro construyó con Redatam: el suelo sin uso no sólo se mide cada diez años por censo: se factura todos los meses. Con una salvedad que corresponde hacer: el baldío tarifario es un \textit{lote sin edificar} con conexión, y la vacancia censal son \textit{viviendas construidas} sin hogar. Los dos padrones miden cosas distintas y no deben sumarse ni tomarse uno por el otro.

\textbf{La tercera es una pregunta abierta, no una afirmación.} Este cuadro rige para el área servida por COSAYSA. La empresa presta servicio en el departamento ---el capítulo \ref{cap:defensas} registra contrataciones suyas con destino a La Caldera y La Calderilla, y el capítulo \ref{cap:amparo} sus captaciones--- e intervino a fojas 68 del expediente de El Durazno; lo que no se sabe es qué parte del ejido cubre. Establecer cuántas conexiones de La Caldera y Vaqueros están categorizadas como baldío --- y cuántas de esas corresponden a los loteos que este libro estudia --- requiere pedirle el padrón a la empresa o al Ente Regulador.

Si ese pedido prospera, el capítulo sobre vacancia deja de apoyarse en una sola fuente decenal y pasa a tener una serie mensual con nombre y dirección. Es, de todos los pedidos pendientes de este libro, el que más cambiaría su capítulo más fuerte.
